\documentclass[11pt,a4paper]{article}
\usepackage[margin=2.2cm]{geometry}
\usepackage{xeCJK}
\usepackage{amsmath,amssymb}
\usepackage{graphicx}
\usepackage{booktabs}
\usepackage{multicol}
\usepackage{tcolorbox}
\usepackage{enumitem}
\usepackage{titlesec}
\usepackage{fancyhdr}
\usepackage{hyperref}
\usepackage{longtable}

\setCJKmainfont{Noto Serif CJK TC}
\setCJKsansfont{Noto Sans CJK TC}

% Style definitions
\titleformat{\section}{\normalfont\Large\bfseries\sffamily}{\thesection}{1em}{}
\titleformat{\subsection}{\normalfont\large\bfseries\sffamily}{\thesubsection}{1em}{}

\pagestyle{fancy}
\fancyhf{}
\fancyhead[L]{\sffamily 國中生物科複習講義}
\fancyhead[R]{\sffamily\leftmark}
\fancyfoot[C]{\thepage}

\newtcolorbox{analysisbox}[1]{%
  colback=blue!5!white,
  colframe=blue!75!black,
  fonttitle=\bfseries\sffamily,
  title={#1},
  arc=1mm
}

\newtcolorbox{learningpoints}[1]{%
  colback=green!5!white,
  colframe=green!60!black,
  fonttitle=\bfseries\sffamily,
  title={#1},
  arc=1mm
}

\newtcolorbox{questionbox}[1]{%
  colback=white,
  colframe=gray!30,
  fonttitle=\bfseries\sffamily,
  coltitle=black,
  title={#1},
  arc=1mm,
  boxrule=0.5pt,
  bottom=1mm,
  top=1mm
}

\begin{document}

\title{\Huge\bfseries\sffamily 國中生物科複習講義}
\author{\large 簡志祥}
\date{2026年}
\maketitle

\tableofcontents
\newpage
\section{生命世界}
\subsection{科學方法}
\begin{learningpoints}{【學習重點整理】}
\begin{enumerate}[label=\arabic*.]
  \item \textbf{【探究歷程】}：觀察 $\rightarrow$ 提出問題 $\rightarrow$ 參考文獻 $\rightarrow$ 提出假設 $\rightarrow$ 設計實驗 $\rightarrow$ 分析實驗結果 $\rightarrow$ 提出結論。如果結論不符合假設，則必須『修正假設』並重新設計實驗。
  \item \textbf{【假設的定義】}：假設是針對觀察到的現象或提出的問題，所做出的『可能解釋』，必須是可以經由實驗驗證的。假設不等於最終結論。
  \item \textbf{【實驗設計原則】}：必須設置『實驗組』與『對照組』進行比較。除了『操縱變因』不同外，其餘『控制變因』必須完全相同。
  \item \textbf{【變因分類】}：
  (1) 『操縱變因』：實驗中唯一改變的變因（例如：探究水分對種子發芽的影響，操縱變因即為『水分』）。
  (2) 『控制變因』：實驗中除操縱變因外，所有其他必須保持相同的變因（例如：溫度、光照、種子數量、容器等），以排除干擾。
  (3) 『應變變因』：實驗所觀測到的結果或數據（例如：發芽率）。
  \item \textbf{【實驗的可重複性】}：單次實驗的偶然誤差較大，因此實驗設計必須使用『足夠的樣本數量』（如每組用50顆種子而非1顆）且必須進行『重複實驗』以求取平均值，從而提高實驗數據的可信度。
  \item \textbf{【科學史與學說】}：假設若經過科學界多次且多方的實驗證實，並得到普遍共識，即可形成『學說』。學說仍有被新科學證據修正或推翻的可能。
  \item \textbf{【會考常見陷阱】}：考題常出現同時改變兩個變因（如同時改變溫度與光照）的錯誤實驗設計，此時實驗結果將無法判斷是由何者引起，該設計無效。
  \item \textbf{【精選考法 - 圖表判讀解題策略】}：會考高頻考點。需仔細識讀橫軸（操縱變因，如時間、溫度）與縱軸（應變變因，如反應速率、體溫）。特別注意曲線的拐點、最高點（最適合值）及趨勢變化（如溫度過高活性降為零的變性特徵，低溫降低但不失活）。
\end{enumerate}
\end{learningpoints}

\subsection{顯微鏡}
\begin{learningpoints}{【學習重點整理】}
\begin{enumerate}[label=\arabic*.]
  \item \textbf{【顯微鏡類型比較】}：
  (1) \textbf{複式顯微鏡}：光線穿透標本，成像為『上下左右相反』的平面影像。適合觀察薄且透光的標本（如細胞、切片）。
  (2) \textbf{解剖顯微鏡}：光線從標本表面反射，成像為『與實物同向且立體』的影像。適合觀察較大、不透光標本（如昆蟲肢體、花朵結構）。
  \item \textbf{【放大倍率與視野變化】}：
  放大倍率 = 目鏡倍率 $\times$ 物鏡倍率。目鏡越長，倍率越低；物鏡越長，倍率越高。
  - 高倍鏡下：視野範圍『變小』、視野亮度『變暗』、細胞數目『變少』、細胞體積『變大』。
  - 低倍鏡下：視野範圍『變大』、視野亮度『變亮』、細胞數目『變多』、細胞體積『變小』。
  \item \textbf{【計算視野中細胞數目的變化】}：
  - 若為『單行排列』的細胞：低倍鏡（100X）換成高倍鏡（400X），放大倍率變為4倍，單行排列的細胞數目變為原來的 $1/4$。
  - 若為『均勻分佈於整個視野』的細胞：放大倍率變為4倍，則視野面積縮小為 $1/16$，視野中能看到的細胞數目變為原來的 $1/16$。
  \item \textbf{【玻片與影像移動法則】}：
  - 在複式顯微鏡下，『像與物移動方向相反』。當影像偏向視野左上方時，欲將影像移至視野中央，應將玻片朝『左上方』移動（簡稱：物隨像走，看到偏哪裡就往哪裡移）。
  - 在解剖顯微鏡下，『像與物移動方向相同』，物像偏左上則玻片往右下移以將其置中。
  \item \textbf{【鏡頭與操作步驟】}：
  操作順序：先使用低倍物鏡尋找焦點（使用粗調節輪），並將目標移至中央，再換高倍物鏡微調焦距（此時『僅能使用細調節輪』，絕不可調整粗調節輪，以免壓碎玻片）。
  \item \textbf{【調光技巧】}：若高倍下視野太暗，可調大光圈或將反光鏡轉為凹面鏡以增加進光量。
  \item \textbf{【氣泡的排除與判斷】}：觀察細胞時如果視野有氣泡（特徵為黑邊、圓形、內部空白），可用鉛筆橡皮擦端輕敲蓋玻片邊緣排除。蓋蓋玻片時應呈45度角慢慢放下，防氣泡產生。
  \item \textbf{【易錯提醒】}：在複式顯微鏡下看到某原生動物游向左上方，要追蹤它，玻片應往左上方移動；在高倍鏡下使用粗調節輪會壓碎玻片並損壞鏡頭。
\end{enumerate}
\end{learningpoints}

\subsection{動植物細胞的比較}
\begin{learningpoints}{【學習重點整理】}
\begin{enumerate}[label=\arabic*.]
  \item \textbf{【細胞構造與生理功能】}：
  (1) 細胞核：含遺傳物質（DNA），為細胞的生命與代謝控制中樞。
  (2) 細胞膜：由雙層脂質構成，為細胞與外界屏障，具『選擇性通透性』，控制物質進出。
  (3) 細胞質：含有多種胞器，為細胞進行多數生化反應（代謝）的場所。
  (4) 粒線體：行『呼吸作用』的主要場所，可分解有機物產生細胞所需的能量（ATP），為細胞的發電機。
  (5) 細胞壁：植物、真菌和細菌特有（主成分為纖維素），具支持與保護作用，維持細胞形狀，但『不具控制物質進出』的功能。
  (6) 葉綠體：含有葉綠素，行『光合作用』製造有機物（葡萄糖）的場所。
  (7) 液泡：儲存水、廢物及營養物質。植物細胞通常有一大液泡，可維持細胞膨壓；動物細胞液泡較小且多。
  \item \textbf{【動植物細胞之主要差異】}：植物細胞多具有細胞壁、大液泡及葉綠體（限綠色組織細胞）；動物細胞無細胞壁、葉綠體，液泡較小。
  \item \textbf{【實驗操作與觀察】}：觀察細胞時常滴加『亞甲藍液』或『碘液』，其目的在於使『細胞核』染色呈深色，利於在顯微鏡下辨識其構造。
  \item \textbf{【葉綠體與綠色細胞的誤區】}：並非所有植物細胞都有葉綠體。洋蔥表皮細胞（透明保護）、植物根部細胞（避光避綠）皆無葉綠體，不可行光合作用；鴨跖草保衛細胞（半月形、成對、圍成氣孔）含有葉綠體，可行光合作用。
  \item \textbf{【常見考法】}：細胞壁不能控制物質進出，負責控制進出的是細胞膜。細胞核受損會導致細胞生理活動失控並逐漸死亡。
  \item \textbf{【精選考法 - 粒線體考點】}：主要考察呼吸作用的場所與能量產生。所有活的真核細胞均含有粒線體，日夜皆行呼吸作用以產生能量（ATP）；原核生物（細菌、藍菌）無粒線體。
  \item \textbf{【精選考法 - 葉綠體考點】}：光合作用的場所。限於含有葉綠素的綠色植物細胞及保衛細胞。動物細胞、真菌細胞、洋蔥表皮細胞、植物根部細胞皆無葉綠體，不可行光合作用。原核生物可行光合作用但無葉綠體。
\end{enumerate}
\end{learningpoints}

\subsection{物質如何交換}
\begin{learningpoints}{【學習重點整理】}
\begin{enumerate}[label=\arabic*.]
  \item \textbf{【擴散與滲透作用】}：
  (1) 擴散作用：物質由高濃度區域向低濃度區域自然散布的現象，不需消耗能量（如：氧氣與二氧化碳的進出）。
  (2) 滲透作用：水分通過選擇性通透膜的擴散。水分會由低濃度溶液（水多）往高濃度溶液（水少）滲透。
  \item \textbf{【細胞的水分平衡（滲透反應）】}：
  (1) 動物細胞（如紅血球）：置於濃鹽水（高張液）會脫水而『萎縮』；置於生理食鹽水（等張液）『形狀不變』；置於蒸餾水（低張液）會『吸水膨脹而破裂（溶血）』。
  (2) 植物細胞：置於濃鹽水會發生『質壁分離』（細胞質流失收縮，與細胞壁分離）；置於蒸餾水會『吸水膨脹，但因有細胞壁保護支持而不會破裂』。
  \item \textbf{【物質進出細胞膜的管道】}：
  (1) 脂質雙層：微小且無極性的分子（如氧氣、二氧化碳）及水分子（部分）可直接穿過。
  (2) 蛋白質通道/載體：水分子（大量）、無機鹽離子、小分子養分（葡萄糖、胺基酸）必須透過膜上特定的蛋白質通道或載體方可進出。
  (3) 大分子物質（澱粉、蛋白質、脂質）：分子過大，絕對無法直接進出細胞膜，必須在消化道內先被分解成單醣（葡萄糖）、胺基酸、甘油與脂肪酸方能被吸收。
  \item \textbf{【主動運輸】}：物質逆濃度梯度（低濃度 $\rightarrow$ 高濃度）運送，需消耗能量，需要細胞膜上特殊載體蛋白的協助（如根部細胞主動吸收無機鹽、小腸上皮細胞主動吸收葡萄糖）。
  \item \textbf{【易錯提醒】}：澱粉、蛋白質等大分子不能通過細胞膜，必須先分解；紅血球在蒸餾水中會破裂，但植物表皮細胞不會，因為植物細胞有細胞壁。
\end{enumerate}
\end{learningpoints}

\subsection{生物體的組成層次}
\begin{learningpoints}{【學習重點整理】}
\begin{enumerate}[label=\arabic*.]
  \item \textbf{【單細胞與多細胞生物】}：
  (1) 單細胞生物：單個細胞即能獨立執行所有的生命現象（如草履蟲、變形蟲、酵母菌、大腸桿菌）。
  (2) 多細胞生物：細胞特化分工，彼此緊密合作（如人、榕樹）。
  \item \textbf{【組成層次比較】}：
  (1) 動物：細胞 $\rightarrow$ 組織 $\rightarrow$ 器官 $\rightarrow$ 器官系統 $\rightarrow$ 個體（具有『器官系統』層次）。
  (2) 植物：細胞 $\rightarrow$ 組織 $\rightarrow$ 器官 $\rightarrow$ 個體（植物『沒有器官系統』層次）。
  \item \textbf{【植物器官分類】}：
  (1) 營養器官：根（吸收水與無機鹽、錨定）、莖（輸送物質、支持）、葉（光合作用、蒸散作用）。
  (2) 繁殖器官：花、果實、種子（參與有性生殖，傳播後代）。
  \item \textbf{【常考實例判定】}：
  - 植物器官：馬鈴薯（塊莖）、地瓜（塊根）、洋蔥（鱗葉）、竹筍（地下莖）；花、果實、種子是器官層次。
  - 動物組織與器官：血液（含有多種血球與血漿）在動物中屬於『組織』層次。心臟、胃、腎臟屬於『器官』層次。
\end{enumerate}
\end{learningpoints}

\newpage
\section{養分}
\subsection{養分}
\begin{learningpoints}{【學習重點整理】}
\begin{enumerate}[label=\arabic*.]
  \item \textbf{【能量營養素】}：醣類（4 kcal/g）、蛋白質（4 kcal/g）、脂質（9 kcal/g）。是生物體產生能量與生長修補的主要來源。醣類是主要的能量來源；蛋白質是構成人體的主要物質；脂質為儲存能量的物質。
  \item \textbf{【非能量營養素】}：水、維生素、礦物質。不產生能量，但對維持正常生理機能與代謝至關重要。缺乏維生素C易得壞血症，缺維生素D易得軟骨症；缺鐵易貧血，缺鈣易骨質疏鬆。
  \item \textbf{【營養素的測定方法】}：
  (1) 澱粉測定：遇碘液呈『藍黑色』（未遇則為黃褐色）。
  (2) 葡萄糖測定：與本氏液混合加熱後呈綠、黃、橙、紅色（依糖量而定，無糖呈淡藍色）。本氏液檢測必須『隔水加熱』以加速反應。
  \item \textbf{【食物所含能量計算】}：利用燃燒釋放的熱量加熱水，吸熱公式為 $H = M \times \Delta T$。燃燒所測得的熱量通常小於實際食物所含熱量，因燃燒過程中部分熱量會散失到空氣中。
  \item \textbf{【核心考點 - 熱量計算】}：食物燃燒釋放能量，可利用水溫上升進行測定，公式為 $H = M \times \Delta T$（熱量 = 水質量 $\times$ 溫度變化量）。由於燃燒過程中部分熱量會散失至空氣或加熱容器，因此實驗測得的熱值通常低於食物實際所含熱量。
  \item \textbf{【核心考點 - 養分檢驗】}：利用特定指示劑檢測養分。澱粉遇碘液呈藍黑色（未遇呈黃褐色）；葡萄糖與本氏液混合並進行『隔水加熱』後，依葡萄糖含量多寡呈現綠、黃、橙、紅色（無糖呈藍色）。
  \item \textbf{【精選考法 - 圖表解讀能力】}：要求學生能從複雜的示意圖（如心臟解剖圖、消化道養分殘留曲線圖）中，辨識各構造的名稱、順序與對應的生理功能。例如，在養分殘留圖中，小腸是澱粉、蛋白質、脂質全部被化學分解與消化完畢的終點器官。
  \item \textbf{【精選考法 - 把表格資料畫成圖】}：要求學生將實驗所得的數值表格，轉換為對應的折線圖或條形圖。轉換時須注意座標軸單位的均勻性、自變數與應變數的座標軸定位（橫軸為自變數/操縱變因，縱軸為應變數/應變變因）。
\end{enumerate}
\end{learningpoints}

\subsection{酵素}
\begin{learningpoints}{【學習重點整理】}
\begin{enumerate}[label=\arabic*.]
  \item \textbf{【酵素的化學本質與催化作用】}：主要由蛋白質組成，作為生物體內的生物催化劑，可降低化學反應活化能，加速代謝反應的進行。反應前後，酵素本身的質量與性質均不改變，可重複使用。
  \item \textbf{【酵素的特性】}：
  (1) \textbf{專一性}：一種酵素只能催化一種或一類化學反應（例如：唾液澱粉酶只能分解澱粉，不能分解蛋白質）。
  (2) \textbf{受溫度與酸鹼值影響}：高溫（一般高於60度）會使酵素的蛋白質結構『永久變性』而失去活性；低溫會降低酵素活性，但不會使酵素永久失活（回溫後活性可恢復）。
  (3) 各種酵素有其最適合運作的酸鹼值（例如：胃蛋白酶適合在強酸下運作，唾液澱粉酶適合在中性下運作）。
  \item \textbf{【常考實驗：唾液澱粉酶對澱粉的分解】}：
  - 操縱變因：溫度（如0度、37度、90度）或酸鹼值（強酸、中性）。
  - 37度中性環境為最適環境，澱粉被分解為麥芽糖，傳統滴加碘液不呈藍黑色，滴加本氏液加熱呈黃橙色。
  - 90度高溫使澱粉酶變性失活，澱粉未分解，滴加碘液呈藍黑色，滴加本氏液加熱呈淡藍色。
  \item \textbf{【精選考法 - 圖表判讀解題策略】}：會考高頻考點。需仔細識讀橫軸（操縱變因，如時間、溫度）與縱軸（應變變因，如反應速率、體溫）。特別注意曲線的拐點、最高點（最適合值）及趨勢變化（如溫度過高活性降為零的變性特徵，低溫降低但不失活）。
\end{enumerate}
\end{learningpoints}

\subsection{光合作用}
\begin{learningpoints}{【學習重點整理】}
\begin{enumerate}[label=\arabic*.]
  \item \textbf{【光合作用的場所與意義】}：主要在葉肉細胞與保衛細胞的『葉綠體』中進行，將光能轉化為化學能儲存在有機物中。
  \item \textbf{【光合作用的原料與產物】}：原料為『水』與『二氧化碳』，產物為『葡萄糖』、『水』與『氧氣』。光合作用的完整反應式為：$6CO_2 + 12H_2O \xrightarrow{\text{光能、葉綠體}} C_6H_{12}O_6 + 6O_2 + 6H_2O$。
  \item \textbf{【光合作用兩大反應步驟】}：
  (1) \textbf{光反應}（在葉綠體的類囊體/基粒上進行）：必須在光照下進行。葉綠素吸收光能，分解水產生『氧氣』（釋放到大氣中）與能量（提供暗反應使用）。
  (2) \textbf{碳反應}（暗反應，在葉綠體基質中進行）：不需光照直接參與。利用光反應提供的能量與大氣中的『二氧化碳』，經過一系列化學反應合成『葡萄糖』與『水』。
  \item \textbf{【光合作用的產物去向】}：葡萄糖可轉化為『澱粉』儲存，或轉化為蛋白質、脂質以供植物生長所需，亦可直接供細胞呼吸作用產生能量（ATP）。
  \item \textbf{【常考實驗：葉片光合作用產物的檢測】}：
  - 步驟與目的：(1) 葉片包裹鋁箔（對照光照變因）。(2) 放入沸水中加熱（使葉細胞死亡，軟化葉片）。(3) 放入『酒精隔水加熱』（溶出葉綠素以利觀察，必須隔水加熱防酒精燃燒）。(4) 漂洗。(5) 滴加碘液（檢測有無澱粉反應，照光區呈藍黑色，遮光區呈黃褐色）。
  \item \textbf{【精選考法 - 圖表判讀解題策略】}：會考高頻考點。需仔細識讀橫軸（操縱變因，如時間、溫度）與縱軸（應變變因，如反應速率、體溫）。特別注意曲線的拐點、最高點（最適合值）及趨勢變化（如溫度過高活性降為零的變性特徵，低溫降低但不失活）。
  \item \textbf{【精選考法 - 實驗設計考點】}：檢驗學生對科學探究變因控制的理解。題目常要求指出某一實驗設計的缺陷（例如：同時改變兩個變因，或缺乏對照組），或要求設計一組對照實驗來驗證某一特定假設（確保單一變因原則）。
  \item \textbf{【精選考法 - 光合作用與pH的關係】}：光合作用會消耗水中的『二氧化碳』（二氧化碳溶於水呈弱酸性）。因此，當水生植物在光照下進行光合作用時，水中的二氧化碳被消耗，會導致水體的酸鹼值（pH值）『上升』（酸性減弱，變為偏鹼性）。此常與指示劑呈色變化結合命題。
\end{enumerate}
\end{learningpoints}

\subsection{消化系統}
\begin{learningpoints}{【學習重點整理】}
\begin{enumerate}[label=\arabic*.]
  \item \textbf{【消化管的組成與順序】}：口腔 $\rightarrow$ 咽 $\rightarrow$ 食道 $\rightarrow$ 胃 $\rightarrow$ 小腸 $\rightarrow$ 大腸 $\rightarrow$ 肛門。消化管的管壁肌肉收縮可蠕動以混合並推進食物。
  \item \textbf{【消化腺與消化液的功能】}：
  (1) 唾腺：分泌唾液（含唾液澱粉酶，初步分解澱粉）。
  (2) 胃腺：分泌胃液（含強鹽酸與胃蛋白酶。強酸可殺菌並輔助胃蛋白酶初步分解蛋白質）。
  (3) 肝臟：分泌『膽汁』儲存於膽囊，藉由膽管注入小腸。\textbf{注意}：膽汁不含任何消化酵素，但可將脂質乳化成小油滴，以增加與酵素接觸面積。
  (4) 胰臟：分泌胰液（含胰澱粉酶、胰蛋白酶、脂肪酶），注入小腸。是人體最主要的消化腺，可消化澱粉、蛋白質與脂質。
  (5) 腸腺：分泌腸液（含多種消化酵素），在小腸進行最後的分解。
  \item \textbf{【消化器官的酸鹼值與功能】}：口腔（中性環境），胃（強酸，pH~2，殺菌並提供胃蛋白酶最佳活性），小腸（弱鹼，提供胰液、腸液、膽汁最佳活性）。
  \item \textbf{【特殊藥劑消化與病症分析】}：
  - 若某藥劑由蛋白質外殼包裹脂質核心，或由脂質外殼包裹蛋白質核心：脂質必須先在『小腸』被膽汁乳化並被胰脂肪酶分解，隨後其內包物質方可釋出並在小腸被消化與吸收。
  - 膽道閉鎖（膽管阻塞）：膽汁無法排出注入小腸，將直接導致脂質（脂肪）無法被乳化，嚴重影響小腸中脂質的化學性消化，使得大分子脂肪難以被脂肪酶分解。
  \item \textbf{【物理性消化與化學性消化】}：
  - 物理性消化：牙齒咀嚼、胃腸蠕動、膽汁乳化、雞吞沙進入雞胗磨碎食物等，不改變養分分子結構，僅增加其與消化酶接觸的表面積。
  - 化學性消化：唾液、胃液、胰液、腸液中的消化酶將澱粉、蛋白質、脂質分解為小分子養分（葡萄糖、胺基酸、甘油與脂肪酸）。
  \item \textbf{【闌尾位置與營養消化】}：闌尾是盲腸的突起，位於右下腹（大腸起始端），發炎時稱為闌尾炎（俗稱盲腸炎）。日常饅頭（澱粉）的消化起於口腔，肉類（蛋白質）起於胃，沙拉油（脂質）起於小腸。
\end{enumerate}
\end{learningpoints}

\newpage
\section{運輸}
\subsection{植物的運輸}
\begin{learningpoints}{【學習重點整理】}
\begin{enumerate}[label=\arabic*.]
  \item \textbf{【維管束的構造與功能】}：
  (1) 木質部（在內側）：運送『水分』與『無機鹽』。運送方向一律為『由下往上』（單向運送）。
  (2) 韌皮部（在外側）：運送『有機養分』（如糖類）。運送方向為『雙向運送』（可由葉往根，或由儲存器官往上運送到生長點）。
  (3) 形成層：位於木質部與韌皮部之間，可進行細胞分裂使莖加粗（雙子葉植物及木本植物具有，單子葉植物無）。
  \item \textbf{【環狀剝皮的生理影響】}：若將木本植物（雙子葉植物）的莖進行環狀剝皮（剝除形成層外側的韌皮部），葉片產生的有機養分無法向下運送到根部，導致根部先飢餓死亡，隨後根部無法吸收水分，整株植物終將因缺水而枯死。
  \item \textbf{【水分上升的動力來源】}：主要動力為『蒸散作用』（水分從葉片氣孔蒸發產生的拉力，拉力可由葉片傳遞至根部），其次為根壓與水分子的毛細現象。
  \item \textbf{【氣孔的開閉機制】}：由兩側的『保衛細胞』控制。
  - 白天或水分充足時：保衛細胞吸水膨脹，因靠近氣孔側之細胞壁較厚不易伸展，導致細胞彎曲，氣孔張開，進行蒸散作用與氣體交換。
  - 夜晚或缺水時：保衛細胞失水萎縮，氣孔關閉，以減少水分流失並停止氣體交換。
  \item \textbf{【細部觀念 - 年輪的形成】}：由雙子葉木本植物莖內的『形成層』受季節氣候影響所產生的木質部細胞差異形成。春夏溫暖多雨，形成層分裂快，木質部細胞大且壁薄，顏色較淡（稱為早材/春材）；秋冬寒冷乾燥，分裂慢，細胞小且壁厚，顏色較深（稱為晚材/秋材）。深淡相間即形成年輪，可推測樹齡與古氣候。
  \item \textbf{【細部觀念 - 維管束的功能與位置】}：維管束是植物體內運輸物質的管道。木質部偏內側，負責運輸水分與無機鹽，方向一律由下往上；韌皮部偏外側，負責運輸有機養分，方向可雙向運送。葉片中的維管束稱為葉脈。
  \item \textbf{【細部觀念 - 維管束綜合性問題】}：單子葉植物維管束呈散生排列，無形成層，莖無法持續加粗（如竹子、玉米）；雙子葉植物維管束呈環狀排列，木質部與韌皮部之間有『形成層』，可持續分裂使莖加粗（如榕樹、樟樹）。
  \item \textbf{【細部觀念 - 樹皮的生理作用】}：樹皮包含死去的表皮、外側栓皮及木本植物莖中『形成層外側的韌皮部』。韌皮部負責運送葉片光合作用產生的有機養分。若樹皮受環狀剝皮，有機養分向下輸送受阻，根部會因飢餓死亡，最後導致整株植物缺水枯死。
  \item \textbf{【精選考法 - 實驗設計考點】}：檢驗學生對科學探究變因控制的理解。題目常要求指出某一實驗設計的缺陷（例如：同時改變兩個變因，或缺乏對照組），或要求設計一組對照實驗來驗證某一特定假設（確保單一變因原則）。
\end{enumerate}
\end{learningpoints}

\subsection{循環系統}
\begin{learningpoints}{【學習重點整理】}
\begin{enumerate}[label=\arabic*.]
  \item \textbf{【血液的組成與功能】}：
  (1) 血漿：主要成分為水，運送葡萄糖、胺基酸、廢物（二氧化碳、尿素）、激素及抗體。
  (2) 紅血球：無細胞核（成熟後），含『血紅素』（富鐵，可在氧氣濃度高處結合氧，氧氣濃度低處釋放氧），呈雙凹圓盤狀，負責運輸氧氣。
  (3) 白血球：有細胞核，體積最大。負責行吞噬作用或產生抗體，防禦病原體。身體感染受傷發炎化膿時，白血球數量會顯著增加。
  (4) 血小板：無細胞核，體積最小。負責促進血液凝固，防止失血過多。
  \item \textbf{【心臟構造與血液流動】}：心臟由心肌構成。分為左心房、左心室、右心房、右心室。左心室壁最厚（收縮力量最強，負責將血液推向全身體循環）。房室之間、心室與動脈之間有『瓣膜』，以防止血液倒流。血液流動方向：靜脈 $\rightarrow$ 心房 $\rightarrow$ 心室 $\rightarrow$ 動脈。
  \item \textbf{【血管類型比較】}：
  (1) 動脈：將血液帶離心臟。管壁最厚、彈性最大、血流速度最快、血壓最高。
  (2) 靜脈：將血液帶回心臟。管壁較薄、彈性小、內有瓣膜（防止血液逆流）。
  (3) 微血管：介於動靜脈之間。管壁最薄（僅一層內皮細胞）、管腔最窄（紅血球需單行通過）、血流速度最慢。是血液與組織細胞進行物質與氣體交換的唯一場所。
  \item \textbf{【兩大循環路徑】}：
  (1) \textbf{體循環}：左心室 $\rightarrow$ 主動脈 $\rightarrow$ 全身微血管 $\rightarrow$ 上下大靜脈 $\rightarrow$ 右心房。流經全身後，動脈血變為靜脈血。
  (2) \textbf{肺循環}：右心室 $\rightarrow$ 肺動脈 $\rightarrow$ 肺部微血管 $\rightarrow$ 肺靜脈 $\rightarrow$ 左心房. 流經肺後，靜脈血變為動脈血。
  - 腦細胞的代謝廢物（二氧化碳）流回心臟：腦部微血管 $\rightarrow$ 頸靜脈 $\rightarrow$ 上大靜脈 $\rightarrow$ 右心房。最先到達『右心房』。
  - 供應心臟肌肉自身所需氧氣與養分的血管為『冠狀動脈』（源自主動脈基部），若阻塞會導致心肌梗塞。
  \item \textbf{【注射藥劑的路徑分析】}：從手臂靜脈注射藥劑，血液隨靜脈回流：手臂靜脈 $\rightarrow$ 上大靜脈 $\rightarrow$ 右心房 $\rightarrow$ 右心室 $\rightarrow$ 肺動脈 $\rightarrow$ 肺部微血管 $\rightarrow$ 肺靜脈 $\rightarrow$ 左心房 $\rightarrow$ 左心室 $\rightarrow$ 主動脈 $\rightarrow$ 全身。藥劑最先進入心臟的『右心房』。
  \item \textbf{【胎盤與臍帶的循環屏障】}：胎生動物的胎盤隔開母體與胎兒的血液循環，兩者血液不直接混合，物質透過「擴散作用」進行交換。此機制可避免母體內大分子病原體直接感染胎兒。
  \item \textbf{【細部觀念 - 心臟的血液供應】}：心臟肌肉自身所需的氧氣與養分由『冠狀動脈』直接供應。冠狀動脈源自主動脈基部，若冠狀動脈發生粥狀硬化或阻塞，會導致心肌細胞缺氧壞死，引起心肌梗塞。
  \item \textbf{【細部觀念 - 血管的類型與特徵】}：動脈管壁厚、彈性大，血流速度最快，負責將血液帶離心臟；靜脈管壁薄、彈性小、管腔大、內有瓣膜（防止血液倒流），負責將血液帶回心臟；微血管僅由一層細胞構成，血流速度最慢，是物質交換的場所。
  \item \textbf{【細部觀念 - 血管的流向】}：體循環流向：心室 $\rightarrow$ 動脈 $\rightarrow$ 微血管 $\rightarrow$ 靜脈 $\rightarrow$ 心房。血液在動脈中流向遠離心臟的方向，在靜脈中流向心臟的方向。微血管介於動靜脈之間，負責物質雙向擴散交換。
\end{enumerate}
\end{learningpoints}

\newpage
\section{協調}
\subsection{神經系統}
\begin{learningpoints}{【學習重點整理】}
\begin{enumerate}[label=\arabic*.]
  \item \textbf{【神經元】}：神經細胞是神經系統的基本單位，包含細胞體（含細胞核，負責維持生命）與突起（傳導訊息）。神經元分為感覺神經元（傳入中樞）與運動神經元（傳出指令到動器）。
  \item \textbf{【中樞神經系統（CNS）的構造與分工】}：
  (1) 大腦：意識、思考、記憶、語言、感覺、運動的核心控制區。大腦左右半球交叉控制身體對側。
  (2) 小腦：維持身體平衡，協調全身肌肉活動（如運動員、特技表演者小腦較發達）。
  (3) 腦幹：生命中樞。控制心搏、呼吸、血壓，以及吞嚥、咳嗽、打噴嚏、唾液分泌等反射。腦幹若壞死即為腦死。
  (4) 脊髓：傳導腦與周圍神經間的訊息。控制頸部以下器官的反射行為（如膝反射、排尿反射、手碰到熱物縮回反射）。
  \item \textbf{【反應路徑分類】}：
  (1) \textbf{意識行為}：路徑必定經過『大腦』皮質。感受器 $\rightarrow$ 感覺神經 $\rightarrow$ 脊髓 $\rightarrow$ 大腦 $\rightarrow$ 脊髓 $\rightarrow$ 運動神經 $\rightarrow$ 動器（肌肉/腺體）。
  (2) \textbf{反射行為}：路徑『不經過大腦』，由脊髓或腦幹直接發出指令。感受器 $\rightarrow$ 感覺神經 $\rightarrow$ 脊髓（或腦幹） $\rightarrow$ 運動神經 $\rightarrow$ 動器。反射行為能快速反應，避開危險，減少機體受傷。
\end{enumerate}
\end{learningpoints}

\subsection{內分泌系統}
\begin{learningpoints}{【學習重點整理】}
\begin{enumerate}[label=\arabic*.]
  \item \textbf{【內分泌腺與激素特點】}：內分泌腺沒有導管（無管腺），分泌的化學物質（稱為『激素』或『荷爾盟』）直接釋放進入微血管，隨血液循環運送到特定目標器官發揮作用。需要量極少，過多或過少皆會引起生理障礙或疾病。
  \item \textbf{【人體主要內分泌腺及其功能】}：
  (1) 腦垂腺：分泌多種激素，其中生長激素過多會引起巨人症，過少引起侏儒症。腦垂腺是內分泌系統的『總指揮』，其分泌的激素可調控甲狀腺、性腺及腎上腺的分泌活動。
  (2) 甲狀腺：分泌甲狀腺素，可促進新陳代謝與生長發育。嬰兒期分泌不足會得呆小症；成年人分泌過多會得甲狀腺機能亢進（代謝過快、體重減輕、心跳加快、眼球突出）。碘是合成甲狀腺素的必需元素。
  (3) 副甲狀腺：分泌副甲狀腺素，可調節體內鈣與磷的平衡，維持神經與肌肉的正常興奮性。分泌不足會抽搐。
  (4) 胰島：位於胰臟內。分泌『胰島素』以降低血糖（促使血糖轉化為肝醣儲存或被細胞利用）；分泌『升糖素』以升高血糖。胰島素分泌不足會導致糖尿病。
  (5) 腎上腺：分泌腎上腺素。在應對緊急狀況（憤怒、恐懼、逃跑）時大量分泌，使心搏加快、血壓上升、呼吸急促、血糖迅速升高，以應付突發狀況。
  (6) 性腺：睪丸（分泌雄性激素）與卵巢（分泌雌性激素），促進第二性徵的發育並維持生殖機能。
  \item \textbf{【核心考點 - 血糖調控】}：人體血糖維持在約 0.1\% 的恆定範圍。飯後血糖上升，刺激胰島分泌『胰島素』，促使細胞吸收血糖或將其轉化為肝醣儲存以降低血糖；運動或飢餓時，刺激胰島分泌『升糖素』及腎上腺分泌『腎上腺素』，促使肝醣分解為葡萄糖釋入血液以升高血糖。
  \item \textbf{【精選考法 - 圖表判讀解題策略】}：會考高頻考點。需仔細識讀橫軸（操縱變因，如時間、溫度）與縱軸（應變變因，如反應速率、體溫）。特別注意曲線的拐點、最高點（最適合值）及趨勢變化（如溫度過高活性降為零的變性特徵，低溫降低但不失活）。
\end{enumerate}
\end{learningpoints}

\subsection{植物感應}
\begin{learningpoints}{【學習重點整理】}
\begin{enumerate}[label=\arabic*.]
  \item \textbf{【植物的感應分類】}：植物沒有神經系統，但可透過化學物質（生長素）的分佈不均或細胞內水分（膨壓）的改變來產生感應。
  \item \textbf{【向性】}：植物受單方向環境刺激（如光照、重力、水分、觸覺）而產生的生長反應。此反應速度緩慢，且為『不可逆的生長現象』。
  - 機制：以向光性為例，單側光照導致背光側的『生長素』濃度較高，促使背光側細胞快速生長，使莖向光彎曲。根部則對生長素敏感度不同，呈現向地性與背光性。
  - 常見向性：莖的向光性與背地性、根的向地性與向濕性、攀緣植物的向觸性。
  \item \textbf{【膨壓運動（局部運動）】}：植物受刺激時，由於局部細胞內水分（膨壓）快速流失或移入，導致細胞體積快速改變而產生的局部快速運動。此運動不涉及細胞生長，為『可逆的短期感應』。
  - 常見膨壓運動：
    (1) 含羞草的『觸震運動』：受觸碰時，葉柄基部細胞的水分快速擴散到細胞間隙，葉片合攏垂下。
    (2) 捕蠅草的『捕蟲運動』：葉片內側敏感毛受觸動時葉片閉合。
    (3) 氣孔的開閉：保衛細胞吸水膨脹，氣孔張開；失水萎縮，氣孔關閉。
    (4) 酢漿草的『睡眠運動』：夜晚水分流失，葉片下垂折疊。
  \item \textbf{【核心考點 - 植物的向性】}：植物受單側光照、重力、水分等單方向刺激，使生長素在背光側或向地側分佈不均，導致兩側細胞生長速度不同而產生的彎曲生長現象。這是不可逆的生長行為，是植物對環境的長期適應性表現。
  \item \textbf{【細部觀念 - 判斷生長原因】}：植物向光性是由於單側光照導致『生長素』在背光側分佈較多，背光側細胞生長速度快於向光側，使莖朝向光源彎曲。根部則因為對生長素敏感度高，高濃度生長素反而抑制生長，呈現背光生長。
  \item \textbf{【細部觀念 - 植物生長現象的原因】}：生長素在植物體內的分佈受重力或光照影響。在莖中，生長素濃度高促進生長，導致背光側或向地側生長快而彎曲；在根中，生長素濃度高反而抑制生長，導致根呈現向地性。
  \item \textbf{【細部觀念 - 預測生長結果】}：受生長素或重力刺激影響，莖會表現向光性與背地性（向上彎曲生長），根表現向地性（向下彎曲生長）。利用生長素在兩側分佈的不均與各器官對生長素敏感度的差異，可精確預測植物幼苗的生長方向。
\end{enumerate}
\end{learningpoints}

\newpage
\section{恆定}
\subsection{體溫恆定}
\begin{learningpoints}{【學習重點整理】}
\begin{enumerate}[label=\arabic*.]
  \item \textbf{【內溫動物與外溫動物】}：
  (1) \textbf{內溫動物}（鳥類、哺乳類）：能利用大腦內的體溫調節中樞與多種生理機制（如肌肉顫抖、出汗、微血管收縮）來產生或散失熱量，維持體溫相對恆定。
  (2) \textbf{外溫動物}（魚類、兩生類、爬行類）：缺乏體溫調節中樞，主要依靠行為方式（如曬太陽、躲入陰涼處、冬眠）來調適體溫。體溫會隨著環境溫度發生劇烈變化。
  \item \textbf{【內溫動物高溫環境下的生理調節（散熱）】}：
  - 皮膚微血管擴張：血流增加，促進輻射散熱。
  - 汗腺分泌增加：汗水蒸發帶走熱量。
  - 減少食慾與活動：降低體內的產熱率。
  \item \textbf{【內溫動物低溫環境下的生理調節（產熱與防散熱）】}：
  - 皮膚微血管收縮：血流減少，減少熱量散失。
  - 食慾增加：補充能量以供產熱。
  - 肌肉顫抖：肌肉快速收縮釋放熱量。
  - 激素分泌（甲狀腺素、腎上腺素）：提高細胞代謝率以增加產熱量。
  \item \textbf{【易錯提醒】}：外溫動物在冬天活動力極低，因為體溫低導致體內酵素活性極差，無法有效進行化學反應，因而需要冬眠；而內溫動物的冬眠主要是為了因應冬天食物短缺，其體溫在冬眠時雖會下降，但仍維持在一定生理限度內。
  \item \textbf{【細部觀念 - 調節體溫的方式】}：內溫動物在低溫下透過血管收縮減少熱量散失，肌肉顫抖與激素分泌（甲狀腺素、腎上腺素）增加產熱；高溫下透過微血管擴張增加散熱，腦幹中樞調控，汗腺分泌汗水蒸發散熱。外溫動物則依賴曬太陽或避光等行為調節。
  \item \textbf{【細部觀念 - 調節體溫的結果】}：內溫動物（鳥類、哺乳類）能將體溫維持在恆定的狹窄生理範圍內，維持酵素的高催化活性與穩定代謝；外溫動物的體溫隨環境溫度劇烈波動，低溫下酵素活性極低而必須冬眠。
\end{enumerate}
\end{learningpoints}

\subsection{氣體恆定}
\begin{learningpoints}{【學習重點整理】}
\begin{enumerate}[label=\arabic*.]
  \item \textbf{【呼吸作用與呼吸運動的區別】}：
  (1) \textbf{呼吸作用}：細胞利用氧氣分解有機養分（葡萄糖）產生能量（ATP）、二氧化碳和水的『生化反應過程』。反應式：$C_6H_{12}O_6 + 6O_2 + 6H_2O \rightarrow 6CO_2 + 12H_2O + \text{能量}$。在細胞的粒線體中進行。
  (2) \textbf{呼吸運動}：藉由肋骨與膈肌（橫膈）的收縮與舒張，改變胸腔體積以完成吸氣與呼氣的『物理運動過程』。
  \item \textbf{【吸氣與呼氣的胸腔變化】}：
  - 吸氣：肋骨往上、往外移，膈肌（橫膈膜）收縮向下移動，胸腔體積變大，肺泡氣壓小於大氣壓，氣體主動流入肺部。
  - 呼氣：肋骨往下、往內移，膈肌舒張向上移動，胸腔體積變小，肺泡氣壓大於大氣壓，氣體被動排出肺部。
  \item \textbf{【血液二氧化碳濃度的調節】}：
  - 血液中的『二氧化碳濃度』是控制呼吸頻率最主要的化學訊號。運動時，細胞呼吸作用增強，血液中二氧化碳濃度上升，刺激『腦幹』的呼吸中樞，發出指令使呼吸運動加快加深，以排除二氧化碳並吸入氧氣。
  \item \textbf{【植物的氣體交換】}：植物體表面無專門的呼吸器官。氣體交換主要透過葉的『氣孔』、莖的『皮孔』以及根部細胞的擴散作用進行。植物日夜皆進行呼吸作用，光照下光合作用大於呼吸作用，淨表現為釋放氧氣；夜晚只進行呼吸作用，淨表現為釋放二氧化碳。
  \item \textbf{【核心考點 - 呼吸作用】}：細胞在粒線體中利用氧氣將葡萄糖等有機物徹底分解，釋放能量（ATP）並產生二氧化碳與水的生化反應過程。是所有活細胞日夜不停進行的基礎代謝，為生命活動提供動力。
  \item \textbf{【核心考點 - 呼吸運動】}：藉由肋骨與膈肌（橫膈膜）的收縮與舒張，改變胸腔體積使胸腔內氣壓與大氣壓產生壓差，完成吸氣（胸腔變大、肺內壓變小）與呼氣（胸腔變小、肺內壓變大）的物理運動，藉此進行肺部與外界的氣體交換。
  \item \textbf{【細部觀念 - 呼吸作用的過程】}：活細胞在粒線體中，利用氧氣將葡萄糖等有機養分分解，釋放出能量（以ATP形式儲存供生命活動使用），並產生二氧化碳和水作為代謝廢物排出。
  \item \textbf{【細部觀念 - 呼吸作用的實驗裝置】}：常利用萌發的種子（呼吸作用旺盛）進行實驗。種子呼吸作用產生二氧化碳，導入『澄清石灰水』中會產生白色混濁（碳酸鈣沉澱），藉此證實細胞呼吸作用釋放二氧化碳。
  \item \textbf{【細部觀念 - 呼吸運動的過程】}：吸氣時，肋骨上舉、膈肌收縮下移，胸腔體積擴大，肺內氣壓低於大氣壓，氣體吸入；呼氣時，肋骨下垂、膈肌舒張上移，胸腔體積縮小，肺內氣壓高於大氣壓，氣體呼出。
  \item \textbf{【細部觀念 - 呼吸運動的調控】}：由大腦下方的『腦幹』自主控制。當人體運動時，細胞呼吸作用增強，血液中二氧化碳濃度上升，刺激腦幹呼吸中樞發出神經衝動，使肋骨與膈肌加速收縮，從而加快加深呼吸頻率。
\end{enumerate}
\end{learningpoints}

\subsection{水分恆定}
\begin{learningpoints}{【學習重點整理】}
\begin{enumerate}[label=\arabic*.]
  \item \textbf{【排泄作用的定義與含氮廢物】}：排泄是指排除細胞代謝產生的代謝廢物（如二氧化碳、尿素、多餘的水與鹽類）的過程。蛋白質代謝會產生有毒的『氨』，生物必須將其排除或轉化：
  (1) 魚類：毒性最強的氨，需大量水稀釋，直接由鰓擴散排入水中。
  (2) 哺乳類與兩生類：氨在『肝臟』轉化為毒性較低的『尿素』，隨血液運送到『腎臟』，過濾形成尿液排出體外。
  (3) 鳥類與昆蟲：將氨轉化為毒性極低、不溶於水的『尿酸』（固態或半固態），與糞便一同排出，可極大地保存體內水分。
  \item \textbf{【人體泌尿系統的構造與運作】}：腎臟（過濾血液，形成尿液） $\rightarrow$ 輸尿管（輸送尿液） $\rightarrow$ 膀胱（暫存尿液） $\rightarrow$ 尿道（排出尿液）。
  \item \textbf{【人體水分恆定的調節】}：
  - 水分過少（如流汗多、飲水少）：血液滲透壓上升，刺激腦沙體分泌抗利尿激素，作用於腎臟增加對水分的重吸收，導致『尿量減少、尿液濃度變濃』，同時刺激大腦產生『口渴感』，促使主動飲水。
  - 水分過多：抗利尿激素減少，腎臟重吸收水分減少，『尿量增加、尿液變淡』，口渴感消失。
  \item \textbf{【植物水分調節】}：植物主要透過『蒸散作用』散失水分。缺水時氣孔關閉以防止進一步流失；有些植物在土壤水分充足且空氣濕度高時，夜晚水分由葉邊緣的排水齒排出，形成『泌液現象』。
  \item \textbf{【核心考點 - 蛋白質的代謝】}：蛋白質在消化道被分解為胺基酸吸收，細胞代謝胺基酸會產生有毒的『氨』。魚類直接由鰓擴散排入水中；哺乳類在肝臟將氨轉化為低毒性的尿素，由腎臟過濾形成尿液排出；鳥類與昆蟲則將其轉化為不溶於水的尿酸隨糞便排出以防水分流失。
  \item \textbf{【精選考法 - 圖表判讀解題策略】}：會考高頻考點。需仔細識讀橫軸（操縱變因，如時間、溫度）與縱軸（應變變因，如反應速率、體溫）。特別注意曲線的拐點、最高點（最適合值）及趨勢變化（如溫度過高活性降為零的變性特徵，低溫降低但不失活）。
\end{enumerate}
\end{learningpoints}

\newpage
\section{生殖}
\subsection{細胞分裂與減數分裂}
\begin{learningpoints}{【學習重點整理】}
\begin{enumerate}[label=\arabic*.]
  \item \textbf{【染色體的基本概念與大小比較】}：染色體位於細胞核內，由DNA與蛋白質組成，攜帶遺傳基因。構造大小關係為：細胞核 > 染色體 > DNA > 基因。基因是DNA分子上的特定遺傳片段。
  \item \textbf{【體細胞與生殖細胞的染色體組成】}：
  - 人類體細胞有23對（46條）染色體。其中22對（44條）為體染色體，1對（2條）為性染色體。男性染色體組成為44+XY，女性為44+XX。
  - 人類生殖細胞（精子或卵子）有23條染色體（單套1n）。男性精子染色體組成為22+X或22+Y，女性卵子一律為22+X。
  \item \textbf{【細胞分裂與減數分裂的詳細比較】}：
  (1) \textbf{細胞分裂}（生長、修補與無性生殖）：染色體複製1次，細胞分裂1次。分裂後同源染色體不分離，產生『2個』子細胞，每個子細胞染色體數目不變（2n $\rightarrow$ 2n）。手受精卵傷口癒合、細胞增生均藉由此方式完成。
  (2) \textbf{減數分裂}（產生配子/有性生殖）：染色體複製1次，細胞連續分裂2次。第一次分裂時『同源染色體分離』；第二次分裂時『複製的染色體單體分離』。產生『4個』子細胞，每個子細胞染色體數目減半（2n $\rightarrow$ 1n）。在女性的卵巢（產生卵子）與男性的睪丸（產生精子）中進行。
  \item \textbf{【分生組織與觀察部位】}：植物的根尖、莖頂之細胞分裂進行旺盛，是觀察細胞分裂的最佳取材部位；成熟的洋蔥表皮細胞、人體成熟紅血球則不進行分裂。
  \item \textbf{【雙胞胎的遺傳組成分析】}：
  - 同卵雙胞胎：由『同一個受精卵』分裂發育而來，兩者的遺傳基因組成100\%完全相同，性別也必然相同。
  - 異卵雙胞胎：由『兩個不同的卵子與兩個精子分別受精』發育而來，遺傳基因不完全相同，性別可同可不同。
  \item \textbf{【體細胞基因的一致性】}：同一個體之不同器官的體細胞皆由受精卵經由『細胞分裂』而來，其『基因組成完全相同』。不同器官外形與功能各異是由於特定基因在不同細胞中選擇性表現的結果。
  \item \textbf{【細部觀念 - 分裂的過程】}：包括染色體複製（遺傳物質加倍）、染色體排列在細胞中央、複製的染色體分離並均勻分配移向細胞兩極、細胞質分裂形成子細胞。細胞分裂用於身體生長、修補與無性生殖。
  \item \textbf{【細部觀念 - 分裂時染色體變化】}：細胞分裂時染色體複製1次、分裂1次，分裂前後染色體數目不變（2n $\rightarrow$ 2n）；減數分裂時染色體複製1次、連續分裂2次，第一次分裂同源染色體分離，第二次分裂染色體單體分離，子細胞染色體數目減半（2n $\rightarrow$ 1n）。
  \item \textbf{【細部觀念 - 染色體在哪裡】}：染色體主要存在於真核細胞的『細胞核』內。在細胞不分裂時呈細絲狀的染色質狀態；當細胞準備分裂時，染色質會高度螺旋化、變粗變短，形成顯微鏡下清晰可見的條狀『染色體』。
  \item \textbf{【細部觀念 - 染色體的判斷】}：同源染色體是指大小、形狀相似，一條來自父親、一條來自母親的成對染色體。在正常體細胞中，染色體成對存在（雙套，2n）；在生殖細胞中，同源染色體分離，僅剩單套染色體（單套，1n）。
  \item \textbf{【細部觀念 - 染色體的組成】}：染色體是由雙螺旋結構的脫氧核糖核酸（DNA）分子和鹼性蛋白質共同纏繞、摺疊組成的複合物，在有絲分裂與減數分裂時進行精確分配。
  \item \textbf{【細部觀念 - 染色體與基因綜合問題】}：基因位於染色體上。同源染色體相對位置上控制同一性狀的基因稱為等位基因（等位基因可相同如TT，或不同如Tt）。減數分裂時，同源染色體分離，其上的等位基因也隨之分離進入不同的配子中。
  \item \textbf{【細部觀念 - 分生組織細胞分裂旺盛】}：植物根尖與莖頂含有『分生組織』，細胞核大、細胞質濃，能進行快速且旺盛的『細胞分裂』使根伸長、莖長高。是顯微鏡下觀察細胞分裂各階段特徵（如染色體排列分離）的最佳實驗取材材料。
  \item \textbf{【細部觀念 - 基因在哪裡】}：基因是控制生物性狀的遺傳片段，位於染色體上的特定位置。染色體由雙螺旋結構的DNA分子和蛋白質共同組成，基因即為DNA分子上的特定核苷酸序列片段。
\end{enumerate}
\end{learningpoints}

\subsection{無性生殖}
\begin{learningpoints}{【學習重點整理】}
\begin{enumerate}[label=\arabic*.]
  \item \textbf{【無性生殖的特裝與優缺點】}：不需經過雌雄配子的結合，僅由母體細胞行細胞分裂直接產生後代。優點是後代能完全保留母體優良性狀，且繁殖速度極快；缺點是後代的遺傳基因與母體完全相同，缺乏變異，當環境劇烈改變時，適應力極低，可能導致全族群滅絕。
  \item \textbf{【無性生殖的常見類型與實例】}：
  (1) 分裂生殖：單細胞生物將自身分裂為兩個子個體（如草履蟲、變形蟲、細菌）。
  (2) 出芽生殖：母體側邊長出芽體，成熟後脫離母體發育（如酵母菌、水螅）。
  (3) 孢子繁殖：利用特殊的孢子萌發成新個體（如黴菌、蕈菇等真菌及蕨類植物的無性階段）。
  (4) 斷裂生殖：身體斷裂成數段，各段再行再生作用發育為新個體（如渦蟲、海星）。
  (5) 營養器官繁殖：利用植物的營養器官進行繁殖。
     - 塊根繁殖：地瓜。
     - 莖繁殖：馬鈴薯、草莓、洋蔥、萬年青與玫瑰。
     - 葉繁殖：落地生根、石蓮花。
  (6) 組織培養：利用植物細胞的全能性，在無菌培養基中大量繁殖同基因植株。
\end{enumerate}
\end{learningpoints}

\subsection{動物有性生殖}
\begin{learningpoints}{【學習重點整理】}
\begin{enumerate}[label=\arabic*.]
  \item \textbf{【受精方式的比較】}：
  (1) \textbf{體外受精}：精子與卵子在親體外的水環境中結合。多見於水生生物（魚類、兩生類、珊瑚）。產卵量極大，但受精率低，胚胎無親代保護。
  (2) \textbf{體內受精}：精子在母體內與卵子結合。多見於陸生生物（爬行類、鳥類、哺乳類、昆蟲）。精子游動有液體介質限制，受精率高，擺脫了對水生環境的依賴，並多具有堅硬蛋殼或在子宮內保護。
  \item \textbf{【胚胎發育方式比較】}：
  (1) \textbf{卵生}：胚胎在體外孵化。發育所需的養分完全由『卵黃』提供。爬行類、鳥類、昆蟲為體內受精卵生；多數魚類與兩生類為體外受精卵生。
  (2) \textbf{胎生}：胚胎在母體子宮內發育。透過『胎盤』與『臍帶』由母體提供氧氣與養分，並排出廢物。多數哺乳類屬於胎生。注意：母體與胎兒的血管並不直接相通，物質是透過擴散作用交換。
  (3) \textbf{卵胎生}：卵在母體內孵化後再產出。胚胎發育所需的養分仍完全由『卵黃』提供，母體僅提供保護作用，不提供養分（如大肚魚、部分毒蛇與鯊魚）。
  \item \textbf{【卵的構造】}：鳥蛋中的蛋白與卵黃是細胞產物，真正的卵細胞僅包含『胚盤、卵黃及卵黃膜』（胚盤含有細胞核）。蛋白和卵黃提供胚胎發育所需的養分與水分，蛋殼具保護與呼吸氣體交換作用。繫帶起固定卵黃的作用。
\end{enumerate}
\end{learningpoints}

\subsection{人類生殖}
\begin{learningpoints}{【學習重點整理】}
\begin{enumerate}[label=\arabic*.]
  \item \textbf{【人類生殖器官與生理功能】}：
  - 男性：主要器官為『睪丸』，負責產生精子與分泌雄性激素（維持第二性徵）。
  - 女性：主要器官為『卵巢』，負責產生卵子與分泌雌性激素、黃體素（維持子宮內膜增厚與懷孕狀態）。
  \item \textbf{【受精與著床發育路徑】}：
  - 卵子排出後進入『輸卵管』。精子由陰道游入，在『輸卵管』中段與卵子相遇並結合完成受精。
  - 受精卵一邊行細胞分裂，一邊朝向子宮移動，最後在『子宮壁』著床，發育為胚胎與胎兒。
  - 懷孕期間，胎兒與母體血液在『胎盤』處進行氣體與養分交換，但兩者血管不相通，血液不直接混合。
  \item \textbf{【月經週期與激素調控】}：子宮內膜隨卵巢排卵週期（約28天）發生增厚與剝落。若卵子未受精，子宮內膜剝落伴隨出血，即為月經。排卵約在兩次月經正中間。
\end{enumerate}
\end{learningpoints}

\subsection{植物有性生殖}
\begin{learningpoints}{【學習重點整理】}
\begin{enumerate}[label=\arabic*.]
  \item \textbf{【花的構造與功能】}：花是被子植物的生殖器官。
  - 花被：花瓣與花萼。
  - 雄蕊：由『花絲』與『花藥』組成。花藥中產生花粉粒，花粉粒萌發產生精細胞。
  - 雌蕊：由『柱頭』、『花柱』及『子房』組成。子房內含有『胚珠』，胚珠內產生卵細胞。
  \item \textbf{【傳粉與受精過程】}：花粉落在雌蕊柱頭上的過程稱為『傳粉』。花粉受柱頭刺激萌發長出『花粉管』，穿過花柱到達子房的胚珠，精細胞沿花粉管運送與胚珠內的卵結合。花粉管的形成使植物受精完全『擺脫對水介質的依賴』，是陸生植物演化的關鍵特徵。
  \item \textbf{【受精後的發育對應關係】}：
  - 受精卵 $\rightarrow$ 胚（新一代植物體）。
  - 胚珠 $\rightarrow$ 種子。
  - 子房壁 $\rightarrow$ 果皮。
  - 子房 $\rightarrow$ 果實。
  一個子房內如果有多個胚珠，受精後果實內就會有多顆種子（例如：西瓜、番茄有多顆種子；桃子、芒果僅有一顆種子，因為子房內只有一個胚珠）。
\end{enumerate}
\end{learningpoints}

\newpage
\section{遺傳}
\subsection{遺傳法則}
\begin{learningpoints}{【學習重點整理】}
\begin{enumerate}[label=\arabic*.]
  \item \textbf{【孟德爾遺傳學的核心概念】}：控制生物性狀的因子稱為『基因』。基因在體細胞中成對存在。染色體上的同源基因可分為顯性（用大寫字母表示，如T）與隱性（用小寫字母表示，如t）。
  \item \textbf{【基因型與表現型】}：
  - 基因型：控制性狀的基因組合（如：TT、Tt、tt）。
  - 表現型：生物外觀顯現出的性狀（如：高莖、矮莖）。
  - 同型合子：成對基因相同（TT, tt）；異型合子：成對基因不同（Tt，僅表現顯性性狀高莖）。
  \item \textbf{【棋盤方格法與比例預測】}：以一對性狀雜交為例（Tt $	imes$ Tt）：
  - 子代基因型比例：TT : Tt : tt = 1 : 2 : 1。
  - 子代表現型比例：顯性性狀 : 隱性性狀 = 3 : 1。
  - 若為 Tt $	imes$ tt（測交）：子代表現型比例為 1 : 1，基因型 Tt : tt = 1 : 1。
\end{enumerate}
\end{learningpoints}

\subsection{人類的遺傳}
\begin{learningpoints}{【學習重點整理】}
\begin{enumerate}[label=\arabic*.]
  \item \textbf{【人類的染色體分類】}：人體細胞共有23對（46條）染色體。
  - 體染色體：前22對，控制一般生理性狀，男女皆同。
  - 性染色體：第23對，決定性別。女性為XX，男性為XY（Y染色體體積小於X）。
  \item \textbf{【人類性別決定機制】}：行減數分裂產生配子時，女性產生的卵細胞皆含X染色體；男性產生的精子半數含X，半數含Y。受精卵發育性別取決於受精精子的種類。
  - 精子(X) + 卵子(X) $\rightarrow$ 女(XX)；精子(Y) + 卵子(X) $\rightarrow$ 男(XY)。生男生女機率各為 50\%，且每一胎皆為獨立事件。
  \item \textbf{【單基因遺傳與ABO血型】}：由一對對位基因控制。人類ABO血型受 $I^A, I^B, i$ 三個等位基因控制（屬於複等位基因遺傳）。
  - A型血基因型：$I^A I^A$ 或 $I^A i$。
  - B型血基因型：$I^B I^B$ 或 $I^B i$。
  - AB型血基因型：$I^A I^B$（共顯性）。
  - O型血基因型：$ii$。
  \item \textbf{【多基因遺傳與性聯遺傳】}：
  - 多基因遺傳：一種性狀受兩對或多對基因共同控制（如：身高、皮膚顏色）。子代表現型呈連續性常態分佈，顯性基因個數越多，特徵越顯著。
  - 性聯遺傳：控制性狀的基因位於性染色體上（多在X染色體）。如紅綠色盲、血友病。男性物要唯一的X染色體帶有致病隱性基因即會發病，故男性患病率遠高於女性。
\end{enumerate}
\end{learningpoints}

\subsection{突變}
\begin{learningpoints}{【學習重點整理】}
\begin{enumerate}[label=\arabic*.]
  \item \textbf{【突變的定義與遺傳規律】}：生物基因或染色體的構造與數量發生異常改變的現象。突變是隨機且多數有害的，但少數有利的突變是生物演化與物種多樣性的原料。
  - \textbf{遺傳法則}：突變若發生在『生殖細胞』，可遺傳給子代；若僅發生在『體細胞』，則不會遺傳給後代。
  \item \textbf{【誘發突變的因素】}：
  (1) 物理因素：紫外線、X光、伽馬射線、放射性物質。
  (2) 化學因素：亞硝酸鹽、黃麴毒素、焦油、塑化劑。
  (3) 生物因素：部分病毒的感染。
  \item \textbf{【人類常見的遺傳性疾病】}：多由隱性基因突變或染色體異常引起。
  - 單基因隱性遺傳病：白化症、紅綠色盲（性聯遺傳）、血友病。
  - 染色體數目異常：唐氏症（第21對體染色體多出一條，共47條）。
  \item \textbf{【近親通婚的危害】}：近親通婚會使隱性致病基因相遇並結合為同型合子的機率大幅增加，提高子代遺傳病的患病率。因此法律禁止近親通婚。
\end{enumerate}
\end{learningpoints}

\subsection{生物科技}
\begin{learningpoints}{【學習重點整理】}
\begin{enumerate}[label=\arabic*.]
  \item \textbf{【基因轉殖技術（基因工程）】}：將一種生物的特定基因（目標基因）切割下來，導入另一種生物的DNA中，使其表現出該基因控制的性狀。
  - 實例：(1) 將人類胰島素基因轉入大腸桿菌，大量生產高純度胰島素。(2) 基因選殖生產抗體或生長激素。(3) 螢光魚與抗蟲作物。
  - 意義：打破物種間生殖隔離的限制，實現跨物種的基因轉移。
  \item \textbf{【複製動物技術】}：利用體細胞核移植技術。將甲動物的體細胞核移入乙動物去核的卵細胞中，發育並植入丙代孕母體產下個體。
  - 遺傳規律：複製後代的遺傳物質完全與提供細胞核的甲動物相同，與乙及丙無關。此技術屬於『無性生殖』。
\end{enumerate}
\end{learningpoints}

\newpage
\section{演化}
\subsection{天擇與人擇}
\begin{learningpoints}{【學習重點整理】}
\begin{enumerate}[label=\arabic*.]
  \item \textbf{【拉馬克的用進廢退說】}：主張生物器官經常使用會發達，不使用會退化，且這種『後天獲得的性狀可以遺傳給後代』。糾正：後天獲得的性狀未改變生殖細胞的DNA，因此『無法遺傳』給後代。
  \item \textbf{【達爾文的自然選擇說（天擇）】}：生物演化最核心的機制。包括四個步驟：
  (1) \textbf{過度繁殖}：生物產生的後代數量遠超過環境資源負荷。
  (2) \textbf{生存競爭}：個體間為爭奪有限資源進行競爭。
  (3) \textbf{個體差異}：同種個體之間存在可遺傳的變異。
  (4) \textbf{適者生存}：環境選擇了具有適應性性狀的個體生存並繁殖後代，不適應者被淘汰。
  天擇決定演化的方向。環境改變，適應的標準也會隨之改變（例如：噴灑殺蟲劑後，具有抗藥性基因的昆蟲比例逐代增加，此為天擇的結果）。
  \item \textbf{【人擇的定義與實例】}：人類根據自身需求與喜好，篩選並繁殖具有特定性狀的生物（如：野生甘藍培育出高麗菜、育種高產量作物）。人擇速度遠快於天擇，但常導致生物遺傳多樣性降低，在野外生存能力下降。
\end{enumerate}
\end{learningpoints}

\subsection{化石}
\begin{learningpoints}{【學習重點整理】}
\begin{enumerate}[label=\arabic*.]
  \item \textbf{【化石的定義與形成】}：古代生物的遺體、遺跡或排泄物被泥沙迅速掩埋，經過漫長的地質作用石化而成。多存在於『沉積岩』中。化石是生物演化最直接的證據。
  \item \textbf{【化石隱含的資訊與地層分布】}：
  - 越古老的地層（越深處）：其中的化石生物構造越簡單、越低等，多為水生無脊椎生物。
  - 越晚近的地層（越淺處）：化石生物構造越複雜、越高等，陸生脊椎生物化石增多。
  - 化石可以推測當時的古氣候、古環境（如珊瑚化石指示溫暖清澈的淺海環境；煤炭代表古代高大蕨類森林）。
  \item \textbf{【標準化石與活化石】}：
  - 標準化石：生存時期短但分佈廣，可用於判定地層年代（如古生代的三葉蟲，中生代的恐龍）。
  - 活化石：結構在漫長地質年代中變化極小，且至今仍存活的生物（如鱟、銀杏、腔棘魚），說明其生存環境非常穩定。
  \item \textbf{【核心考點 - 化石與地層關係】}：化石多存於沉積岩中。地層沉積通常越底層越古老，其中的化石生物構造越簡單低等（多為水生無脊椎生物）；越表層地層越晚近，化石生物越複雜高等（多為陸生脊椎生物）。由化石可推測古氣候與古環境（如溫暖淺海的珊瑚化石，潮濕陸地的煤炭地層）。
\end{enumerate}
\end{learningpoints}

\subsection{地質年代}
\begin{learningpoints}{【學習重點整理】}
\begin{enumerate}[label=\arabic*.]
  \item \textbf{【地質年代的劃分依據】}：科學家根據地層中的主要標準化石種類及地質事件（如大滅絕、火山爆發），將地球歷史劃分為前寒武紀、古生代、中生代和新生代。
  \item \textbf{【各世代代表性生物演化歷程】}：
  (1) \textbf{前寒武紀}：生命起源（海洋中出現單細胞原核生物），隨後出現真核生物及多細胞無脊椎動物。
  (2) \textbf{古生代}：海洋無脊椎動物（三葉蟲）繁盛，魚類興起。隨後陸地出現維管束植物。兩生類與昆蟲興起，脊椎動物開始登陸。
  (3) \textbf{中生代}：『爬行動物』（恐龍）與裸子植物繁盛。鳥類與哺乳類在此時期開始出現。中生代末期（6500萬年前）小行星撞擊地球導致恐龍滅絕。
  (4) \textbf{新生代}：『哺乳動物』與被子植物繁盛。人類在此時期興起並演化。
  \item \textbf{【核心考點 - 地質年代的事件】}：地質年代劃分為古生代、中生代與新生代。古生代以三葉蟲為代表，後期蕨類繁榮、魚類與兩生類興起；中生代以恐龍（爬行類）與裸子植物為代表；新生代以哺乳類、被子植物（開花植物）與人類繁榮為代表。各代末期常伴隨大規模生物大滅絕。
\end{enumerate}
\end{learningpoints}

\subsection{演化樹}
\begin{learningpoints}{【學習重點整理】}
\begin{enumerate}[label=\arabic*.]
  \item \textbf{【演化樹的原理與閱讀方法】}：演化樹是用來表示生物間親緣關係與演化歷程的圖表。
  - 節點：代表共同祖先。越晚近的分支點，表示親緣關係越近。
  - 分支距離：在演化樹上，如果兩種生物具有越晚近的共同祖先，說明兩者的親緣關係越密切。
  \item \textbf{【同源構造與分子證據】}：
  - 同源構造: 來源相同、基本結構相似，但功能可能發生適應性改變的器官（如蝙蝠的翼手、鯨魚的胸鰭、人類的前肢）。同源構造證明生物具有共同祖先。
  - DNA及蛋白質序列相似度：DNA核苷酸序列或胺基酸序列越相似，生物間親緣關係越近。這是現代系統分類學最客觀的證據。
  \item \textbf{【核心考點 - 判斷演化關係】}：利用演化樹（樹狀圖）的分支點來判斷。分支點代表共同祖先，離現在越近的分支點表示兩物種的親緣關係越近；若兩者在演化樹上的分開路徑越短，說明兩者的DNA或蛋白質序列相似度越高，親緣關係越親近。
  \item \textbf{【細部觀念 - 共同祖先】}：在演化樹中，分支點代表生物的共同祖先。分支點越晚近，表示兩者分開演化的時間越短，親緣關係越近；例如，鳥類與爬蟲類具有較晚近的共同祖先，說明兩者的親緣關係比鳥類與兩生類更為密切。
\end{enumerate}
\end{learningpoints}

\newpage
\section{分類}
\subsection{命名與分類}
\begin{learningpoints}{【學習重點整理】}
\begin{enumerate}[label=\arabic*.]
  \item \textbf{【林奈雙名法學名規則】}：學名 = 屬名 + 種小名（斜體字或加底線）。
  (1) 屬名：屬名相同，表示親緣關係極近，同屬不同物種。字首大寫。
  (2) 種小名：形容詞，字首小寫，描述特徵或產地。
  - 實例：犬的學名為 \textit{Canis familiaris}，狼為 \textit{Canis lupus}。屬名相同，親緣關係近。
  \item \textbf{【分類階層】}：界 $\rightarrow$ 門 $\rightarrow$ 綱 $\rightarrow$ 目 $\rightarrow$ 科 $\rightarrow$ 屬 $\rightarrow$ 種。階層越高，包含的生物種類越多，共同特徵越少，親緣關係越遠；階層越低，包含的生物種類越少，共同特徵越多，親緣關係越近。共同分類階層越低，親緣關係越近。
  \item \textbf{【『種』的生物學定義】}：分類學的最基本單位。在自然狀態下能互相交配，並能產生『具有生育能力』的後代。若能交配但後代無生育能力，則親代非同種生物（如馬和驢生下的騾無生育能力）。
\end{enumerate}
\end{learningpoints}

\subsection{病毒與原核生物}
\begin{learningpoints}{【學習重點整理】}
\begin{enumerate}[label=\arabic*.]
  \item \textbf{【病毒的特徵】}：
  - 結構：無細胞結構。僅由外層的『蛋白質外殼』與內部的『遺傳物質』（DNA或RNA）組成。
  - 生理：在活細胞外無任何生命現象；必須寄生在宿主活細胞內方能進行增殖。不屬於五界系統中的任何一界。抗生素對病毒無效。
  \item \textbf{【原核生物界特徵與分類】}：
  - 特徵：單細胞。具有細胞壁與細胞膜，但『無核膜包被的細胞核』，遺傳物質散在細胞質中。細胞內『無具膜的胞器』（如無粒線體、葉綠體）。
  - 分類實例：
    (1) 細菌：絕大多數行異營生活（如大腸桿菌、乳酸菌）。
    (2) 藍菌（藍綠藻）：『可行光合作用』自營製造養分（注意：藍菌無葉綠體，其光合色素散在細胞質膜上）。如念珠藻。
  \item \textbf{【精選考法 - 給生物列特徵】}：提供某一種特定生物（如蝙蝠、水筆仔、落地生根），要求指出其分類歸屬及生理特徵組合（如：蝙蝠具毛髮、胎生、哺乳、內溫；水筆仔具維管束、胎生苗繁殖、被子植物）。
\end{enumerate}
\end{learningpoints}

\subsection{真菌界}
\begin{learningpoints}{【學習重點整理】}
\begin{enumerate}[label=\arabic*.]
  \item \textbf{【真菌界的特徵】}：
  - 結構：屬於真核生物，具有由幾丁質構成的『細胞壁』。
  - 生理：『無葉綠體』，絕對無法進行光合作用。行異營生活，主要以『腐生』或『寄生』方式生存。是重要的分解者。
  \item \textbf{【真菌的分類與實例】}：
  (1) 單細胞真菌：『酵母菌』。行出芽生殖，用於釀酒及麵包發酵。
  (2) 多細胞真菌：由『菌絲』組成。
     - 黴菌：如黑黴菌、青黴素。
     - 蕈菇類：香菇、木耳、靈芝。其地上肉眼可見構造為『子實體』，負責產生孢子進行繁殖。
\end{enumerate}
\end{learningpoints}

\subsection{植物界}
\begin{learningpoints}{【學習重點整理】}
\begin{enumerate}[label=\arabic*.]
  \item \textbf{【植物界的共同特徵】}：多細胞真核生物。細胞壁主成分為纖維素。含有葉綠體，行光合作用自營生活。具世代交替現象。
  \item \textbf{【植物界四大類群的演化特徵】}：
  (1) \textbf{蘚苔植物}：『無維管束』，無真正的根、莖、葉結構。矮小，水分靠擴散運送。受精仍需水為介質（如地錢、土馬騌）。
  (2) \textbf{維管束無子植物（蕨類）}：『有維管束』，有真正的根、莖、葉。利用『孢子』進行無性繁殖。受精仍需水為介質（如鐵線蕨、筆筒樹）。
  (3) \textbf{裸子植物}：有維管束。產生種子，但『無花與果實』構造，種子裸露在『毬果』的鱗片上，胚珠外露。常見植物有松、杉、柏、銀杏。
  (4) \textbf{被子植物}：有維管束。具有『花』與『果實』，種子受到果實保護。依子葉數目分為：
     - 單子葉植物：子葉1片、平行脈、維管束散生、鬚根系、花瓣數為3的倍數（如稻、麥、玉米、竹）。
     - 雙子葉植物：子葉2片、網狀脈、維管束環狀（有形成層可加粗）、直根系、花瓣數為4或5的倍數（如榕樹、玫瑰、大豆）。
  \item \textbf{【細部觀念 - 裸子植物的特徵】}：具有維管束，產生種子繁殖，但『無花與果實』構造，胚珠與種子裸露在『毬果』的鱗片上。多以風力傳粉，無雙重受精現象。如松、杉、柏、銀杏。
  \item \textbf{【細部觀念 - 蘚苔和蕨類的差異】}：蕨類『有維管束』，有真正的根、莖、葉結構，利用葉背孢子囊產生的孢子繁殖；蘚苔植物『無維管束』，無真正根莖葉（僅假根），水分靠擴散運送，植株矮小。兩者受精過程皆需要水為介質。
  \item \textbf{【精選考法 - 給生物列特徵】}：提供某一種特定生物（如蝙蝠、水筆仔、落地生根），要求指出其分類歸屬及生理特徵組合（如：蝙蝠具毛髮、胎生、哺乳、內溫；水筆仔具維管束、胎生苗繁殖、被子植物）。
  \item \textbf{【精選考法 - 多種生物的比較】}：考題常將多種不同界/門的生物（如水母、昆蟲、海豚、鯊魚）列出，要求比較其受精方式（體內/體外）、發育方式（卵生/胎生）、體溫調節（內溫/外溫）或呼吸器官（鰓/肺/皮膚）。
\end{enumerate}
\end{learningpoints}

\subsection{動物界}
\begin{learningpoints}{【學習重點整理】}
\begin{enumerate}[label=\arabic*.]
  \item \textbf{【無脊椎動物主要門類特徵】}：
  (1) 刺絲胞動物門：身體輻射對稱，口周圍有觸手（含刺絲胞）。如水母、珊瑚、海葵。
  (2) 扁形動物門：身體扁平，左右對稱，單一開口。如渦蟲、絛蟲。
  (3) 軟體動物門：身體柔軟，多具貝殼，有外套膜. 如蝸牛、烏賊、章魚、蛤蜊。
  (4) 環節動物門：身體細長，有重複的『體節』，無附肢。如蚯蚓、水蛭。
  (5) 節肢動物門：體表有『外骨骼』，附肢分節，需定期『蛻皮』。昆蟲（3對足）、蜘蛛（4對足）、甲殼類（蝦蟹）屬於此門。
  (6) 棘皮動物門：體表有棘，具管足，輻射對稱。如海星、海膽、海參。
  \item \textbf{【脊椎動物亞門五大類群特徵】}：
  (1) 魚類：體表具鱗片。用鰓呼吸。體外受精、卵生（多數）。外溫。
  (2) 兩生類：幼體水生鰓呼吸，成體陸生肺及皮膚呼吸。發育經變態。體外受精、卵生。外溫。如青蛙、山椒魚。
  (3) 爬行類：體表有角質鱗片或骨板，卵具革質硬殼。體內受精、卵生。外溫。如蛇、蜥蜴、烏龜。
  (4) 鳥類：覆羽。前肢特化為翼。骨骼中空。卵具鈣質硬殼。體內受精、卵生。內溫。
  (5) 哺乳類：毛髮。具乳腺，以乳汁哺育幼兒。體內受精，多數為胎生（僅鴨嘴獸、針鼴為卵生）。內溫。
  \item \textbf{【核心考點 - 脊椎動物的特徵】}：體內具有脊椎骨組成的脊柱。包括魚類（鰓呼吸、鰭運動、外溫）、兩生類（幼體水生鰓呼吸，成體陸生肺及皮膚呼吸、外溫、發育需經變態）、爬行類（角質鱗片防失水、陸上產卵、外溫）、鳥類（羽毛、中空骨骼、恆溫、卵生）、哺乳類（毛髮、乳腺育幼、恆溫、多胎生）。
  \item \textbf{【細部觀念 - 棘皮動物的特徵】}：屬於水生無脊椎動物，身體呈輻射對稱（成體），體表有棘狀突起，具有特有的管足系統，利用『管足』進行運動、捕食與氣體交換。如海星、海膽、海參。
  \item \textbf{【細部觀念 - 兩生類的特徵】}：屬於外溫動物，幼體（蝌蚪）用鰓呼吸、在水中生活，成體用肺及潮濕的皮膚（輔助呼吸）在陸地生活。發育需經『變態』過程。受精與產卵必須在水中進行，卵無硬殼保護，無法完全擺脫水生環境。
  \item \textbf{【細部觀念 - 兩生類與爬蟲類的差異】}：爬行類體表覆蓋角質鱗片或骨板，可有效防止水分散失，卵具革質硬殼保護，行體內受精，完全適應陸地生活；兩生類皮膚濕潤無覆蓋物，極易散失水分，卵無硬殼，多行體外受精且必須在水中發育。
  \item \textbf{【細部觀念 - 海豚的分類與哺乳類的特徵】}：海豚生活在海洋中，外形似魚，但體內具肺（用肺呼吸）、體溫恆定、胎生、具乳腺以乳汁哺育幼兒，因此分類上屬於『哺乳綱』而非魚綱。
  \item \textbf{【精選考法 - 多種生物的比較】}：考題常將多種不同界/門的生物（如水母、昆蟲、海豚、鯊魚）列出，要求比較其受精方式（體內/體外）、發育方式（卵生/胎生）、體溫調節（內溫/外溫）或呼吸器官（鰓/肺/皮膚）。
  \item \textbf{【精選考法 - 給生物列特徵】}：提供某一種特定生物（如蝙蝠、水筆仔、落地生根），要求指出其分類歸屬及生理特徵組合（如：蝙蝠具毛髮、胎生、哺乳、內溫；水筆仔具維管束、胎生苗繁殖、被子植物）。
\end{enumerate}
\end{learningpoints}

\newpage
\section{生態}
\subsection{族群與群集}
\begin{learningpoints}{【學習重點整理】}
\begin{enumerate}[label=\arabic*.]
  \item \textbf{【族群的定義與計算】}：
  - 族群：特定時間、生活在同一區域的『同種生物個體』的總和。
  - 變動公式：族群數量變化 = (出生 + 遷入) - (死亡 + 遷出)。
  - 標記重捕法公式：總個體數 / 第一次標記數 = 第二次捕獲總數 / 第二次捕獲中的標記數。應進行多次抽樣並求平均值。
  \item \textbf{【負荷量】}：一個環境受空間、食物、天敵限制，所能維持該族群個體數目的最大極限。族群數量通常會在負荷量上下波動。
  \item \textbf{【群集的定義】}：特定時間、生活在同一區域內『所有不同生物族群』的集合。群集內生物間有複雜的交互關係。
  \item \textbf{【精選考法 - 數據判讀方法】}：考題常給出實驗數據表格（如不同溫度下的發芽率、不同血球計數等）。解題時需先找出對照組與實驗組，觀察數值的變化規律，並排除偶然誤差，以找出操縱變因與應變變因之間的科學因果關係。
\end{enumerate}
\end{learningpoints}

\subsection{交互關係}
\begin{learningpoints}{【學習重點整理】}
\begin{enumerate}[label=\arabic*.]
  \item \textbf{【生物間的交互作用分類】}：
  (1) \textbf{競爭}（-,-）：兩種生物爭奪相同的有限資源（如陽光、食物、空間）。競爭結果常導致一方被排擠。如作物與雜草。
  (2) \textbf{掠食}（+,-）：捕食者殺死並吃掉獵物。如貓頭鷹捕食田鼠。獵物常演化出保護色防禦。
  (3) \textbf{寄生}（+,-）：寄生蟲生活在宿主的體內或體表吸取營養，不立即致死。如跳蚤與狗、菟絲子與綠色植物。
  (4) \textbf{互利共生}（+,+）：雙方皆能獲得利益。如地衣（真菌與綠藻）、小丑魚與海葵、根瘤菌與豆科植物。
  (5) \textbf{片利共生}（+,0）：一方獲利，另一方無利無害。如蘭花附生在大樹上、印魚吸附在鯊魚腹部。
  \item \textbf{【細部觀念 - 生物競爭的原因】}：同種或不同種生物個體，為了爭奪相同的有限生存資源（如食物、陽光、水分、空間、配偶）而發生資源搶奪的關係。資源越匱乏，競爭越劇烈。
  \item \textbf{【細部觀念 - 判斷發生競爭的生物】}：在同一個生態系中，如果兩種生物具有相似的棲地、食物來源或生存空間，它們之間就會發生資源搶奪的『競爭』關係（如：農田中的農作物與雜草爭奪陽光與水分，森林中的獅子與豹爭奪獵物）。
\end{enumerate}
\end{learningpoints}

\subsection{能量的流動}
\begin{learningpoints}{【學習重點整理】}
\begin{enumerate}[label=\arabic*.]
  \item \textbf{【生態系的角色分工與食物鏈】}：
  (1) 生產者：自營生物（綠色植物、藻類、藍菌）。是一切能量的門戶。
  (2) 消費者：異營生物（各級消費者）。
  (3) 分解者：將生物遺體分解為無機物釋回環境，以維持物質循環（如細菌、真菌）。
  \item \textbf{【能量流動的單向性與遞減規律】}：
  - 太陽能是生態系能量的根本來源。能量在食物鏈中呈『單向流動，無法循環』。每次能量轉移時，僅有約 10\% 的能量流向下一營養級（十分之一定律），其餘 90\% 主要以『熱能散失』。
  - 食物鏈通常不超過五級，因為能量在傳遞中快速遞減。
  \item \textbf{【生物累積作用】}：環境中不易分解且不易排出的毒性物質（如重金屬汞、鉛，以及戴奧辛、DDT），沿著食物鏈傳遞時，會逐級累積，使得『營養級越高的生物體內累積的毒物濃度越高』。頂端捕食者（老鷹、人類）受害最深。
  \item \textbf{【細部觀念 - 從能量流動關係判斷生態系角色】}：在生態系中，生產者（植物）能吸收太陽能並進行光合作用，是一切能量的來源；消費者（動物）通過捕食獲得有機物；分解者（細菌、真菌）分解動植物遺體，將有機物轉化為無機物釋回環境，以維持物質循環。
  \item \textbf{【細部觀念 - 能量塔的分析】}：能量在食物鏈傳遞中，每往上一層僅有約 10\% 轉移，其餘 90\% 以熱能散失或被生物自身消耗。因此，營養級越高，獲得的能量越少，呈現下寬上窄的金字塔形。
  \item \textbf{【細部觀念 - 從能量塔判斷生物】}：能量在食物鏈傳遞中逐級遞減（十分之一定律）。因此在能量金字塔中，位於底層的生物（生產者）所含能量與個體數量通常最多；越往頂層（高階消費者），所含能量與個體數量越少，呈現下寬上窄的金字塔形。
\end{enumerate}
\end{learningpoints}

\subsection{物質的循環}
\begin{learningpoints}{【學習重點整理】}
\begin{enumerate}[label=\arabic*.]
  \item \textbf{【碳循環】}：碳是構成有機物的重要元素。
  - 大氣中二氧化碳透過生產者的『光合作用』進入生物體合成有機物。
  - 生物有機碳透過生物的『呼吸作用』、分解者的『分解作用』或化石燃料的『燃燒』，再次轉化為二氧化碳釋回大氣中，形成循環。
  \item \textbf{【氮循環】}：氮是合成蛋白質與核酸的必需元素。
  - 固氮作用：大氣中的氮氣必須透過『固氮細菌』（如根瘤菌）或雷電，轉化為植物可吸收的銨鹽或硝酸鹽。
  - 植物吸收氮化合物合成有機氮，經食物鏈傳遞。
  - 脱氮作用：遺體經分解者分解產生氨，再由硝化細菌與脫氮細菌作用，將氮元素釋放回大氣中。
  \item \textbf{【水循環】}：地球上的水分透過太陽加熱蒸發為水蒸氣，以及植物的『蒸散作用』、動物的呼吸與排泄散失到大氣中。再藉由降雨、降雪等降水形式回到陸地與海洋，供生物吸收利用。水循環是物質循環中物理狀態改變最顯著的循環。
  \item \textbf{【核心考點 - 氮循環】}：氮氣占大氣78\%，但植物無法直接吸收。必須靠固氮細菌（如豆科植物根瘤菌）進行固氮作用轉化為銨鹽或硝酸鹽，方可被植物吸收用以合成蛋白質與核酸。動植物遺體由分解者分解產生氨，再經硝化與脫氮細菌流回大氣。
\end{enumerate}
\end{learningpoints}

\subsection{生態}
\begin{learningpoints}{【學習重點整理】}
\begin{enumerate}[label=\arabic*.]
  \item \textbf{【生態系的類型與分佈決定變因】}：主要受『雨量』與『溫度』（緯度與海拔）決定。
  \item \textbf{【陸域生態系】}：
  (1) 森林生態系：雨量充沛。植物以高大喬木為主，森林層次複雜。生物多樣性最高。
  (2) 草原生態系：年雨量中等，有明顯乾濕季。植物以草本植物為主。動物多具奔跑能力。
  (3) 沙漠生態系：年雨量極稀少，日夜溫差極大。植物具發達之防禦與節水構造（如仙人掌針狀葉）。
  \item \textbf{【水域生態系】}：
  (1) 淡水生態系：分為流動水體（溪流，氧氣充足，生物具流線型身體或吸盤）與靜止水體（湖泊、池塘）。
  (2) 河口生態系：淡水與海水交界處。鹽度與水分波動劇烈，營養鹽極豐富。生產力極高，生物多具耐鹽、耐旱特徵（如紅樹林、招潮蟹、彈塗魚）。單一物種數量多。
  (3) 海洋生態系：
     - 近海區：光照充足、養分豐富，生物種類繁多。
     - 遠洋區：表層有光照，浮游植物為主要生產者；深海區無光、水溫低、壓力大，生產者缺乏，深海生物多以有機碎屑（海雪）或鯨落為食。
\end{enumerate}
\end{learningpoints}

\subsection{生物多樣性}
\begin{learningpoints}{【學習重點整理】}
\begin{enumerate}[label=\arabic*.]
  \item \textbf{【生物多樣性的三個層次】}：
  (1) \textbf{遺傳多樣性}：同種生物個體之間基因組合的差異（例如：人類的不同血型與膚色、狗的多種品種、野生稻與農作稻的抗病基因差異）。遺傳多樣性越高，當環境發生變遷時，族群中存活下適應個體的機率越高，族群越不易滅絕。
  (2) \textbf{物種多樣性}：一個生態系內『物種數量的多寡』以及『各物種個體分佈的均勻度』。物種多樣性越高，食物網越複雜，生態系越穩定。
  (3) \textbf{生態系多樣性}：一個區域內生態系類型的多樣性。生態系多樣性高能提供多樣的棲地與環境功能，維持地球生態圈的運作平衡。
  \item \textbf{【保護生物多樣性的意義與措施】}：生物多樣性提供人類食物、醫藥原料、維持氣候恆定、水土保持及休閒價值。主要保護措施包括建立自然保護區、野生動物保育法、國際公約，以及保護棲地（就地保護是成效最佳的方式）。
  \item \textbf{【細部觀念 - 外來種入侵】}：外來物種引進本土後，由於『缺乏天然天敵』的控制，且具有極強的適應力與繁殖力，大量繁殖並搶奪本土物物生存資源（食物、棲地），造成本土物物瀕危或滅絕，嚴重降低當地的生物多樣性。
  \item \textbf{【細部觀念 - 生物多樣性的觀念】}：包含三個層次：遺傳（基因）多樣性（個體間基因差異）、物種多樣性（物種多寡與均勻度）及生態系多樣性（生存環境多樣化）。三個層次相輔相成，共同維持地球生態系統的穩定與健康。
  \item \textbf{【細部觀念 - 推測生物多樣性高的棲地】}：氣候溫暖、雨量豐沛、環境複雜且受到人為干擾少的棲地（如熱帶雨林、珊瑚礁生態系），能提供多樣化的棲地與充足的食物資源，其物種多樣性與生態系多樣性最高。
  \item \textbf{【細部觀念 - 維維持生物多樣性的方法】}：建立自然保護區（保護棲地，就地保護是最有效的方法）、制定保育法律（如野生動物保育法）、控制污染與外來種引進、加強大眾環境教育。
  \item \textbf{【細部觀念 - 維持生物多樣性的目的】}：維持生態系的穩定與平衡，提供人類醫藥、食物、原料的來源，保持基因庫的完整性以利作物改良，維護大自然自淨能力與水土保持。
\end{enumerate}
\end{learningpoints}

\subsection{環境}
\begin{learningpoints}{【學習重點整理】}
\begin{enumerate}[label=\arabic*.]
  \item \textbf{【溫室效應與全球暖化】}：人類大量燃燒化石燃料，釋放過量『二氧化碳』及甲烷等溫室氣體，導致大氣吸收地表反射的紅外線輻射增加，全球氣溫上升。
  - 影響：兩極冰川融化、海平面上升、極端氣候頻發、全球生物分布帶往高緯度或高海拔遷移。
  \item \textbf{【臭氧層的破壞】}：人類過去使用『氟氯碳化物』（CFCs）釋放到平流層，分解臭氧。導致臭氧層變薄甚至出現空洞。
  - 影響：地表紫外線（UV）輻射增加，導致皮膚癌、白內障率上升，並破壞海洋浮游生物，危害食物鏈基礎。
  \item \textbf{【外來種入侵的生態威脅】}：引進外來生物（如福壽螺、小花蔓澤蘭、綠鬣蜥、巴西鐵樹、非洲大蝸牛）。由於在本土『缺乏天敵控制』，且適應力與繁殖力極強，大肆繁殖。
  - 威脅：搶奪本土物種的生存空間與食物，造成本土物種瀕危或滅絕，嚴重降低本土生物多樣性。
  \item \textbf{【優養化現象】}：含有大量氮、磷等營養鹽的廢水排入靜止水體，導致藻類爆發性繁殖。藻類大量死亡後，分解者（細菌）大量繁殖並消耗水中溶解氧，導致魚蝦因缺氧窒息大量死亡，水體發臭崩潰。
  \item \textbf{【精選考法 - 數據判讀方法】}：考題常給出實驗數據表格（如不同溫度下的發芽率、不同血球計數等）。解題時需先找出對照組與實驗組，觀察數值的變化規律，並排除偶然誤差，以找出操縱變因與應變變因之間的科學因果關係。
  \item \textbf{【精選考法 - 優養化發生的原因】}：考題常涉及環境保護。主要成因為家庭廢水、化學肥料或洗滌劑中的氮、磷等營養鹽大量排入湖泊或池塘等靜止水體，導致水體營養過剩。
  \item \textbf{【精選考法 - 優養化發生的結果】}：水體優養化引發藻類在短時間內爆發性繁殖，覆蓋水面阻擋光照。藻類死亡後，分解者（細菌）大量繁殖，過程中消耗水中大量溶解氧，導致水生動物（魚蝦）因缺氧窒息死亡，水體發臭崩潰。
\end{enumerate}
\end{learningpoints}

\newpage
\section{其他}
\subsection{其他綜合觀念}
\begin{learningpoints}{【學習重點整理】}
\begin{enumerate}[label=\arabic*.]
  \item \textbf{【跨章節綜合分析方法】}：解答綜合性生物試題時，應注意將『生理作用』、『系統調控』與『生態意義』相連結。例如，運動時人體呼吸加快（氣體恆定），心搏加快（運輸系統），此過程受腦幹（神經系統）與腎上腺素（內分泌系統）的共同調節。
  \item \textbf{【生物學史重要里程碑與科學家貢獻】}：
  - 虎克：利用自製顯微鏡觀察軟木塞切片，首次發現並命名細胞（Cell，實為細胞壁）。
  - 林奈：創立雙名法，奠定生物分類學的基石。
  - 達爾文：提出自然選擇說，出版《物種起源》，確立現代演化論。
  - 孟德爾：以豌豆進行雜交實驗，發現遺傳因子，被尊稱為遺傳學之父。
  \item \textbf{【日常生活與環境中的生物學應用】}：
  - 生物防治：利用天敵控制害蟲數量（如用寄生蜂控制介殼蟲），以減少農藥使用，保護環境。
  - 基因檢測：利用分子生物技術檢測遺傳缺陷或疾病風險。
  - 減碳與環保：認識個人生活中的碳足跡，落實資源回收、減少使用一次性塑膠，以減緩全球暖化與環境破壞。
\end{enumerate}
\end{learningpoints}

\newpage
\end{document}
